%%%%%%%%%%%%%%%%%%%%%%%%%%%%%%%%%%%%%%%%%%%%%%%%%%%%%%%%%%%%%%%%%%%%%%
% Lớp 8 | Chương III. Tứ giác | Hình bình hành
% Dạng 8H3?12-2: Dấu hiệu nhận biết hình bình hành
%%%%%%%%%%%%%%%%%%%%%%%%%%%%%%%%%%%%%%%%%%%%%%%%%%%%%%%%%%%%%%%%%%%%%%


% nguồn: De_AI_Toan8_cuoi_hoc_ki_i_1 (nhập 21/11/2026)
\begin{bt}%[8H3H12-2]%[ma:fec90cb1]
	Cho hình bình hành $ABCD$. Gọi $E$ và $F$ lần lượt là trung điểm của các cạnh $AB$ và $CD$.
a) Chứng minh tứ giác $AECF$ là hình bình hành.
b) Chứng minh $AF \parallel CE$ và $AF = CE$.
	\loigiai{
		a) Vì $ABCD$ là hình bình hành nên $AB \parallel CD$ và $AB = CD$. \\
Vì $E$, $F$ lần lượt là trung điểm của $AB$, $CD$ nên $AE = \dfrac{1}{2}AB$ và $CF = \dfrac{1}{2}CD$, suy ra $AE = CF$. \\
Lại có $AE \parallel CF$ (do $AB \parallel CD$). \\
Xét tứ giác $AECF$ có $AE \parallel CF$ và $AE = CF$, nên $AECF$ là hình bình hành. \\
b) Vì $AECF$ là hình bình hành nên hai cạnh đối $AF$ và $CE$ song song và bằng nhau, tức $AF \parallel CE$ và $AF = CE$.
	}
\end{bt}

% nguồn: De_AI_Toan8_cuoi_hoc_ki_i_1 (nhập 21/11/2026)
\begin{bt}%[8H3H12-2]%[ma:0e7dbc9f]
	Cho hình bình hành $ABCD$ có hai đường chéo $AC$ và $BD$ cắt nhau tại $O$. Trên đường chéo $AC$ lấy hai điểm $M$ và $N$ sao cho $AM = CN$ (với $M$ nằm giữa $A$ và $O$, $N$ nằm giữa $C$ và $O$). Chứng minh tứ giác $BMDN$ là hình bình hành.
	\loigiai{
		Vì $ABCD$ là hình bình hành có hai đường chéo cắt nhau tại $O$ nên $O$ là trung điểm của $AC$ và $BD$, suy ra $OA = OC$ và $OB = OD$. \\
Vì $M$ nằm giữa $A$ và $O$, $N$ nằm giữa $C$ và $O$ nên $OM = OA - AM$ và $ON = OC - CN$. \\
Mà $OA = OC$ và $AM = CN$ nên $OM = ON$, tức $O$ là trung điểm của $MN$. \\
Xét tứ giác $BMDN$ có hai đường chéo $BD$ và $MN$ cắt nhau tại trung điểm $O$ của mỗi đường, nên $BMDN$ là hình bình hành.
	}
\end{bt}

% nguồn: De_AI_Toan8_cuoi_hoc_ki_i_1 (nhập 21/11/2026)
\begin{bt}%[8H3H12-2]%[ma:06f37cd8]
	Cho hình bình hành $ABCD$. Kẻ $AE \perp BD$ tại $E$ và $CF \perp BD$ tại $F$.
a) Chứng minh tam giác $ABE$ bằng tam giác $CDF$.
b) Chứng minh tứ giác $AECF$ là hình bình hành.
	\loigiai{
		a) Vì $ABCD$ là hình bình hành nên $AB = CD$ và $AB \parallel CD$, suy ra $\widehat{ABE} = \widehat{CDF}$ (so le trong). \\
Xét hai tam giác vuông $ABE$ và $CDF$ ($\widehat{AEB} = \widehat{CFD} = 90^\circ$) có $AB = CD$ và $\widehat{ABE} = \widehat{CDF}$. \\
Suy ra $\triangle ABE = \triangle CDF$ (cạnh huyền – góc nhọn). \\
b) Từ $\triangle ABE = \triangle CDF$ suy ra $AE = CF$. \\
Mặt khác $AE \parallel CF$ (vì cùng vuông góc với $BD$). \\
Xét tứ giác $AECF$ có $AE \parallel CF$ và $AE = CF$, nên $AECF$ là hình bình hành.
	}
\end{bt}

% nguồn: De_AI_Toan8_cuoi_hoc_ki_i_1 (nhập 21/11/2026)
\begin{bt}%[8H3H12-2]%[ma:5ddada24]
	Cho tam giác $ABC$ có đường trung tuyến $AM$. Trên tia đối của tia $MA$ lấy điểm $D$ sao cho $MD = MA$.
a) Chứng minh tứ giác $ABDC$ là hình bình hành.
b) Biết $\widehat{BAC} = 65^\circ$, tính số đo $\widehat{BDC}$ và $\widehat{ABD}$.
	\loigiai{
		a) Vì $AM$ là đường trung tuyến của tam giác $ABC$ nên $M$ là trung điểm của $BC$. \\
Theo cách lấy điểm $D$, $M$ cũng là trung điểm của $AD$. \\
Xét tứ giác $ABDC$ có hai đường chéo $AD$ và $BC$ cắt nhau tại trung điểm $M$ của mỗi đường, nên $ABDC$ là hình bình hành. \\
b) Vì $ABDC$ là hình bình hành nên $\widehat{BDC} = \widehat{BAC} = 65^\circ$ (hai góc đối bằng nhau) và $\widehat{ABD} = 180^\circ - \widehat{BAC} = 180^\circ - 65^\circ = 115^\circ$.
	}
\end{bt}

% nguồn: chuyen_dang (nhập 22/11/2026)
\begin{ex}%[8H3N12-2]%[ma:e55a8c52]
	Trong các  mệnh đề sau, mệnh đề nào là mệnh đề đảo của mệnh đề “Nếu tứ giác $ABCD$ là hình bình hành thì tứ giác $ABCD$ có các góc đối bằng nhau”?
	\choice 
		{Tứ giác $ABCD$ là hình bình hành suy ra tứ giác $ABCD$  có các góc đối bằng nhau}
		{\True Nếu tứ giác $ABCD$ có các góc đối bằng nhau  thì tứ giác $ABCD$ là hình bình hành}
		{Tứ giác $ABCD$ là  hình bình hành nên tứ giác $ABCD$ có các góc đối bằng nhau}
		{Nếu tứ giác $ABCD$ là hình bình hành thì tứ giác $ABCD$ có các góc đối bằng nhau} 
	 \loigiai{
	Mệnh đề đảo của mệnh đề $P\Rightarrow Q$ là mệnh đề $Q\Rightarrow  P$.}
\end{ex}

% nguồn: chuyen_dang (nhập 22/11/2026)
\begin{bt}%[8H3H12-2]%[ma:06dd9f89]
	Trong mặt phẳng toạ độ $Oxy$, cho tứ giác $ABCD$ có $A(x_A;y_A)$, $B(x_B;y_B)$, $C(x_C;y_C)$, $D(x_D;y_D)$. Chứng minh tứ giác $ABCD$ là hình bình hành khi và chỉ khi $x_A+x_C=x_B+x_D$ và $y_A+y_C=y_B+y_D$.
	\loigiai{
		Tứ giác $ABCD$ là hình bình hành khi và chỉ khi
		$$
		\overrightarrow{AB}=\overrightarrow{DC} \Leftrightarrow \heva{&x_B-x_A=x_C-x_D\\&y_B-y_A=y_C-y_D} \Leftrightarrow \heva{&x_B+x_D=x_C+x_A\\&y_B+y_D=y_C+y_A.}
		$$
	}
\end{bt}
